\mxunit{automatismes-ch08}{Variations et extrémums d'une fonction}

\begin{fluencysheet}{jeu.~11~févr.~2027}{Croissance, décroissance, monotonie}
\mxfq{Échauffement · Flash}{IMG}{$f(x)=-2x+5$. Calculer $f(1)$ et $f(3)$\,; que constate-t-on\,?}{$f(1)=3$, $f(3)=-1$~: l'ordre des images est inverse de celui des abscisses ($f$ décroît).}{}
\mxfq{Cœur · QCM}{MOD}{$f$ est croissante sur $I$ signifie que, pour tous $a<b$ dans $I$~:\par\smallskip\noindent A.~$a\leqslant b$ seulement\quad B.~$f(a)\geqslant f(b)$\par\noindent C.~$f(a)\leqslant f(b)$\quad D.~$f(a)=f(b)$}{C.\ $f(a)\leqslant f(b)$}{Croissante~: conserve l'ordre\,; décroissante~: le renverse.}
\mxfq{Cœur + · Graphique}{REG}{Voici la courbe de $f(x)=x^2$ sur $\intff{-2}{2}$~:\par\smallskip\begin{tikzpicture}[scale=0.45,>=stealth]
\draw[very thin,gray!45] (-2,-1) grid (2,4);
\draw[->] (-2,0)--(2.3,0); \draw[->] (0,-1)--(0,4.3);
\foreach \i in {-2,-1,1,2}\draw (\i,.1)--(\i,-.1) node[below,font=\tiny]{\i};
\foreach \j in {-1,1,2,3,4}\draw (.1,\j)--(-.1,\j) node[left,font=\tiny]{\j};
\draw[thick,mxink,domain=-2:2,samples=60,smooth] plot (\x,{\x*\x});

\end{tikzpicture}\par Donner les variations de $f$.}{Décroissante sur $\intff{-2}{0}$, croissante sur $\intff{0}{2}$.}{Expression utilisée pour le tracé~: $f(x)=x^2$.}
\mxfq{Retour · Flash}{PRO}{Quel coefficient multiplicateur correspond à une hausse de $\pc{15}$\,? À une baisse de $\pc{15}$\,?}{$\num{1,15}$\,; $\num{0,85}$.}{}
\mxfq{Fil rouge · Flash}{DIV}{$457$ est-il divisible par $3$\,? Est-il divisible par $9$\,? Justifier.}{Somme des chiffres $=16$, non multiple de $3$~: ni par $3$, ni par $9$.}{}
\end{fluencysheet}

\begin{fluencysheet}{lun.~1\textsuperscript{er}~mars~2027}{Tableau de variations\,; maximum et minimum sur un intervalle}
\mxfq{Échauffement · Flash}{IMG}{Calculer $x^2$ pour $x=-2$, $0$, $3$.}{$4$\,; $0$\,; $9$.}{}
\mxfq{Cœur · Méthode}{REG}{$f$ est décroissante sur $\intff{-3}{1}$ avec $f(-3)=5$, $f(1)=-1$, puis croissante sur $\intff{1}{4}$ avec $f(4)=4$. Donner le maximum et le minimum de $f$ sur $\intff{-3}{4}$.}{Maximum $5$ atteint en $x=-3$ ($5>4$)\,; minimum $-1$ atteint en $x=1$.}{Comparer les valeurs aux extrémités de l'intervalle ET aux changements de variation.}
\mxfq{Cœur + · Vrai/faux}{REG}{Vrai ou faux\,? Justifier. Le maximum d'une fonction sur $\intff{a}{b}$ est toujours atteint en $a$ ou en $b$.}{Faux. Contre-exemple~: $f(x)=-x^2+4$ sur $\intff{-2}{2}$~: maximum $4$ atteint en $0$.}{}
\mxfq{Retour · Tableau}{PNP}{Ce tableau est-il de proportionnalité\,?\par\smallskip\mxarraystretch{1.15}\begin{tabular}{|c|c|c|c|}\hline
$x$ & 1 & 2 & 3 \\ \hline
$y$ & 5 & 9 & 13 \\ \hline
\end{tabular}\par }{Non~: $\dfrac51\neq\dfrac92$ (accroissement constant $+4$ mais $y(0)\neq0$).}{}
\mxfq{Fil rouge · Méthode}{GEO}{Le triangle de côtés $5$, $12$ et $13$ est-il rectangle\,? Quelle propriété faut-il utiliser\,?}{$5^2+12^2=169=13^2$~: réciproque du théorème de Pythagore, rectangle (angle droit opposé à $13$).}{}
\end{fluencysheet}

\begin{fluencysheet}{mar.~2~mars~2027}{Fonctions affines~: taux d'accroissement et variations}
\mxfq{Échauffement · Flash}{IMG}{$f(x)=3x-1$. Calculer $\dfrac{f(4)-f(1)}{4-1}$.}{$\dfrac{11-2}{3}=3$~: c'est la pente $m$.}{}
\mxfq{Cœur · Flash}{MOD}{Quel est le sens de variation de $f(x)=-4x+7$\,?}{$m=-4<0$~: $f$ est strictement décroissante sur $\R$.}{}
\mxfq{Cœur + · Démonstration}{LIT}{Montrer que $f(x)=5x-2$ est croissante sur $\R$.}{Pour $a<b$~: $f(b)-f(a)=5(b-a)>0$, donc $f(a)<f(b)$.}{Modèle de rédaction pour la démonstration au programme.}
\mxfq{Retour · Tableau}{LEC}{Voici les âges des élèves d'un club.\par\smallskip\mxarraystretch{1.15}\begin{tabular}{|c|c|c|c|c|}\hline
Âge & $10$ & $11$ & $12$ & $13$ \\ \hline
Effectif & 3 & 6 & 10 & 6 \\ \hline
\end{tabular}\par Combien d'élèves ont au plus $11$ ans\,? Quelle est la fréquence de ces élèves\,?}{$3+6=9$ élèves\,; $\dfrac{9}{25}=\pc{36}$.}{}
\mxfq{Fil rouge · Programme → formule}{REG}{Un programme de calcul consiste à multiplier $x$ par $3$, puis à soustraire $2$. Donner $f(x)$ puis $f(4)$.}{$f(x)=3x-2$\,; $f(4)=10$.}{}
\end{fluencysheet}

\begin{fluencysheet}{jeu.~4~mars~2027}{Fonctions valeur absolue et carré~: signe et variations}
\mxfq{Échauffement · Flash}{CAL}{Calculer $|{-}5|$, $(-5)^2$ et $|2-7|$.}{$5$\,; $25$\,; $5$.}{}
\mxfq{Cœur · Méthode}{REG}{Donner le tableau de variations de $x\mapsto|x|$ et son signe.}{Décroissante sur $\intof{-\infty}{0}$, croissante sur $\intfo{0}{+\infty}$\,; minimum $0$ en $0$\,; toujours positive.}{}
\mxfq{Cœur + · Réflexion}{CMP}{Comparer $(-3)^2$ et $(-2)^2$ puis $3^2$ et $2^2$. La fonction carré est-elle croissante sur $\R$\,?}{$9>4$ dans les deux cas, mais $-3<-2$~: le carré renverse l'ordre sur $\intof{-\infty}{0}$\,; non, pas croissante sur $\R$.}{}
\mxfq{Retour · Méthode}{PRB}{On lance deux dés équilibrés. Quelle est la probabilité que la somme vaille $7$\,?}{$6$ cas favorables sur $36$~: $\dfrac16$.}{}
\mxfq{Fil rouge · Méthode}{ALG}{Un programme de calcul consiste à choisir un nombre, à ajouter $1$, à élever le résultat au carré, puis à soustraire le carré du nombre de départ. Quel résultat obtient-on pour $6$\,? Quelle expression générale obtient-on, en notant $x$ le nombre de départ\,?}{Pour $6$~: $49-36=13$. En général $(x+1)^2-x^2=2x+1$.}{}
\end{fluencysheet}

\begin{fluencysheet}{lun.~8~mars~2027}{Fonction inverse~: variations\,; position relative de $y=x$ et $y=x^2$}
\mxfq{Échauffement · Flash}{CAL}{Calculer $\dfrac12$, $\dfrac1{-4}$ et $\dfrac1{\num{0,1}}$.}{$\num{0,5}$\,; $-\num{0,25}$\,; $10$.}{}
\mxfq{Cœur · Méthode}{CMP}{Comparer $\dfrac13$ et $\dfrac15$\,; que dit-on des variations de la fonction inverse sur $\intoo{0}{+\infty}$\,?}{$\dfrac13>\dfrac15$ alors que $3<5$~: l'inverse est décroissante sur $\intoo{0}{+\infty}$ (et sur $\intoo{-\infty}{0}$, mais pas sur $\R^*$).}{}
\mxfq{Cœur + · Vrai/faux}{CMP}{Vrai ou faux\,? Justifier. On affirme que, pour tout réel $x>0$, $x^2\geqslant x$.}{Faux. Contre-exemple~: $x=\dfrac12$, $x^2=\dfrac14<\dfrac12$. En fait $x^2\geqslant x$ si et seulement si $x\geqslant1$ (pour $x>0$).}{}
\mxfq{Retour · Flash}{PER}{Quelle est la vue de face et la vue de dessus d'un cylindre posé sur sa base\,?}{Face~: un rectangle\,; dessus~: un disque.}{}
\mxfq{Fil rouge · Prouver}{LIT}{Montrer que le produit de deux entiers pairs est un multiple de $4$.}{Avec $a=2p$, $b=2q$~: $ab=4pq$, multiple de $4$.}{}
\end{fluencysheet}

\begin{fluencysheet}{mar.~9~mars~2027}{Comparer $f(a)$ et $f(b)$\,; problèmes d'optimisation}
\mxfq{Échauffement · Flash}{CMP}{$f(x)=x^2$. Comparer $f(-3)$ et $f(2)$.}{$f(-3)=9>4=f(2)$.}{}
\mxfq{Cœur · Méthode}{REG}{$f(x)=x^2$, $a=-\num{1,5}$ et $b=-2$. Comparer $f(a)$ et $f(b)$ sans calcul, en utilisant les variations.}{$-2<-\num{1,5}\leqslant0$ et $f$ décroissante sur $\intof{-\infty}{0}$~: $f(-2)>f(-\num{1,5})$. Contrôle~: $4>\num{2,25}$.}{}
\mxfq{Cœur + · Optimisation}{MOD}{On considère $A(x)=x(10-x)$ pour $x\in\intff{0}{10}$ (aire d'un rectangle de périmètre $20$). Montrer que $A(x)=25-(x-5)^2$ et en déduire le maximum.}{$25-(x-5)^2=25-x^2+10x-25=x(10-x)$\,; maximum $25$ atteint pour $x=5$ (carré de côté $5$).}{}
\mxfq{Retour · Flash}{STA}{$18$ élèves sur $24$ ont plus de $10$ de moyenne. Donner la fréquence en pourcentage.}{$\dfrac{18}{24}=\pc{75}$.}{}
\mxfq{Fil rouge · QCM}{EQ}{La solution de l'équation $\dfrac{x}{3}+1=\dfrac{7}{2}$ est~:\mxchoices{$\dfrac{15}{2}$}{$\dfrac{9}{2}$}{$\dfrac{21}{2}$}{$\dfrac{5}{6}$}}{A.\ $\dfrac{15}{2}$}{Piège~: oublier de soustraire le $1$ avant de multiplier par $3$ ($\frac{21}{2}$).}
\end{fluencysheet}
